\documentclass[11pt,a4paper]{article}
\usepackage[utf8]{inputenc}
\usepackage[margin=1in]{geometry}
\usepackage{hyperref}
\usepackage{listings}
\usepackage{xcolor}
\usepackage{booktabs}
\usepackage{titlesec}
\usepackage{amsmath,amssymb}

% Package setup
\hypersetup{
    colorlinks=true,
    linkcolor=blue,
    filecolor=magenta,      
    urlcolor=cyan,
    pdftitle={DataKit Official User Guide},
    pdfpagemode=FullScreen,
}

% Code Listing Configuration
\definecolor{codegreen}{rgb}{0,0.6,0}
\definecolor{codegray}{rgb}{0.5,0.5,0.5}
\definecolor{codepurple}{rgb}{0.58,0,0.82}
\definecolor{backcolour}{rgb}{0.96,0.96,0.96}

\lstdefinestyle{pythonstyle}{
    backgroundcolor=\color{backcolour},   
    commentstyle=\color{codegreen},
    keywordstyle=\color{magenta},
    numberstyle=\tiny\color{codegray},
    stringstyle=\color{codepurple},
    basicstyle=\ttfamily\small,
    breakatwhitespace=false,         
    breaklines=true,                 
    captionpos=b,                    
    keepspaces=true,                 
    numbers=left,                    
    numbersep=5pt,                  
    showspaces=false,                
    showstringspaces=false,
    showtabs=false,                  
    tabsize=4
}
\lstset{style=pythonstyle}

\title{\textbf{DataKit: Safety-First Python Data Science Abstraction Layer}\\[0.5em] \large Official Technical User Guide \& Reference Manual}
\author{\textbf{Rajtech69}}
\date{Version 0.1.0 --- August 2026}

\begin{document}

\maketitle

\begin{abstract}
DataKit is a human-oriented, safety-first Python library built directly on top of NumPy, Pandas, Matplotlib, and Seaborn. It provides explicit safety guards against silent broadcasting bugs, non-destructive data cleaning with mandatory confirmation gates, pure Object-Oriented visual rendering without global state-bleed, scikit-learn pipeline dataset preparation, and zero-dependency multi-format report generation. This document provides a complete technical specification and user manual for DataKit v0.1.0.
\end{abstract}

\tableofcontents
\newpage

\section{Introduction \& Design Architecture}

DataKit is designed around five core architectural pillars:
\begin{enumerate}
    \item \textbf{Explicit Safety over Silent Convenience}: Destructive cleaning operations require explicit \texttt{confirm=True} arguments. Silent errors and implicit rank-broadcasting are flagged or prevented.
    \item \textbf{Immutable Pipeline Ergonomics}: Functions that clean or transform datasets return a new \texttt{DataKit} instance containing a modified copy, leaving the original dataset untouched.
    \item \textbf{Zero Matplotlib State-Bleed}: All visualization methods operate purely on Matplotlib \texttt{Figure} and \texttt{Axes} objects, allowing seamless integration into custom subplot grids.
    \item \textbf{Three-Part Explained Exceptions}: All error messages adhere to a three-part structure: \textit{what happened}, \textit{why it matters}, and \textit{suggested fix}.
    \item \textbf{Thin Abstraction Layer}: DataKit orchestrates NumPy and Pandas without introducing redundant custom data structures.
\end{enumerate}

\section{Installation \& Requirements}

DataKit requires Python 3.10 or higher.

\subsection{Core Dependencies}
\begin{itemize}
    \item \texttt{numpy $\ge$ 1.24}
    \item \texttt{pandas $\ge$ 2.0}
    \item \texttt{matplotlib $\ge$ 3.7}
    \item \texttt{seaborn $\ge$ 0.13}
\end{itemize}

\subsection{Installation Options}
\begin{lstlisting}[language=bash]
# Standard Installation via Git
pip install git+https://github.com/Rajtech69/DataKit.git

# Installation with optional ML, Excel, and Arrow extras
pip install "git+https://github.com/Rajtech69/DataKit.git#egg=datakit[all]"
\end{lstlisting}

\section{Core DataKit Class}

The \texttt{DataKit} class in \texttt{datakit.core.datakit} serves as the primary dataset wrapper.

\begin{lstlisting}[language=Python]
import datakit as dk

# Instantiation from CSV, Excel, Parquet, JSON, or DataFrame
data = dk.DataKit("dataset.csv")

# Escape Hatch: Access live DataFrame reference
df = data.df
\end{lstlisting}

\section{Broadcasting \& Safety Module}

The \texttt{dk.safe} module prevents silent NumPy and Pandas shape broadcasting errors.

\begin{lstlisting}[language=Python]
import datakit as dk
import numpy as np

a = np.arange(5)          # Shape (5,)
b = np.arange(5)[:, None]  # Shape (5, 1)

# Issues ImplicitBroadcastWarning naming (5,) - (5, 1) -> (5, 5)
res = dk.safe.subtract(a, b)

# Inspect shapes pre-flight
shape_info = dk.safe.check_shapes(a, b)
print(shape_info.explanation)
\end{lstlisting}

\section{Data Quality \& Non-Destructive Cleaning}

DataKit segregates data inspection from modification.

\subsection{Data Audit}
\begin{lstlisting}[language=Python]
audit = data.audit()
print(audit.summary)
df_issues = audit.to_dataframe()
\end{lstlisting}

\subsection{Non-Destructive Cleaning}
\begin{lstlisting}[language=Python]
# Report-first default (no modifications made)
data.clean()

# Confirmed non-destructive clean
cleaned = data.clean(missing="impute_median", duplicates="drop", confirm=True)
print(cleaned.last_clean_report.diff_summary())
\end{lstlisting}

\section{Object-Oriented Visualization}

All plots return a \texttt{PlotResult} holding \texttt{.fig} and \texttt{.ax} attributes.

\begin{lstlisting}[language=Python]
import matplotlib.pyplot as plt

fig, axes = plt.subplots(1, 2, figsize=(10, 4))
cleaned.plot.hist("age", ax=axes[0])
cleaned.plot.scatter("age", "charges", hue="smoker", trend=True, ax=axes[1])
\end{lstlisting}

\section{Machine Learning Dataset Preparation}

\begin{lstlisting}[language=Python]
ml_data = cleaned.prepare(
    target="charges",
    task="regression",
    scale=True,
    encode="onehot"
)

X_train = ml_data.X_train
pipeline = ml_data.preprocessing_pipeline
\end{lstlisting}

\section{Multi-Format Synthesis Reports}

\begin{lstlisting}[language=Python]
# Export 8-section synthesis report in HTML, Markdown, or JSON
cleaned.report(format="html", path="synthesis_report.html")
cleaned.report(format="markdown", path="report.md")
cleaned.report(format="json", path="report.json")
\end{lstlisting}

\end{document}
